\chapter{Anti-Mullerian Hormone (AMH)}

\section*{What Is This Test?}
This test measures the level of anti-Mullerian hormone (AMH) in your blood. AMH is a glycoprotein hormone produced by the granulosa cells of small, growing ovarian follicles (preantral and small antral follicles less than 4-6 mm in diameter). AMH levels reflect the size of the remaining ovarian follicle pool --- in other words, your ovarian reserve. Unlike FSH or estradiol, AMH levels remain relatively stable throughout the menstrual cycle, making it a convenient test that can be drawn on any day. AMH is used in fertility evaluation to predict ovarian response to stimulation protocols, to diagnose polycystic ovary syndrome (PCOS), and to assess menopausal transition. AMH is also produced by Sertoli cells in the testes; in males, it is used to evaluate cryptorchidism and disorders of sex development in infants, and to assess Sertoli cell function.

\section*{Alternative Names}
AMH; Anti-Mullerian Hormone; Mullerian Inhibiting Substance; MIS; Mullerian Inhibiting Hormone

\section*{Principle}
AMH is measured using a quantitative two-site immunoassay. The most widely used platform is the Beckman Coulter Access/DxI AMH Gen II assay, which uses a sandwich immunoassay format with two monoclonal antibodies: a capture antibody immobilized on magnetic particles and a detection antibody conjugated to alkaline phosphatase. The patient serum is incubated with the capture antibody-coated particles, washed, then incubated with the detection antibody. Chemiluminescent signal is measured and is proportional to AMH concentration. The assay range is typically 0.08 to 22.4 ng/mL (0.57 to 160 pmol/L). An automated extraction step using a pre-treatment reagent was added in the Gen II assay to reduce interference from heterophilic antibodies. Results are reported in ng/mL (nanograms per milliliter) or pmol/L (picomoles per liter). Conversion: 1 ng/mL = 7.14 pmol/L.

\section*{History}
AMH was first identified in the 1940s by Alfred Jost as the substance responsible for regression of the Mullerian ducts during male fetal development \cite{jost1947recherches}. It was initially called Mullerian Inhibiting Substance (MIS). The gene was cloned in 1986 by Cate and colleagues, enabling development of immunoassays \cite{cate1986isolation}. Clinical application in female fertility assessment began in the early 2000s when studies demonstrated that AMH correlated with antral follicle count (AFC) and ovarian response to gonadotropin stimulation. The first commercial AMH assay (Diagnostic Systems Laboratories) became available in 2002. The Beckman Coulter Gen II assay was introduced in 2009 and became the most widely used platform. The addition of an extraction step in 2013 (Gen II with extraction) reduced false-positive results from heterophilic antibodies. Current guidelines from the American Society for Reproductive Medicine (ASRM) and European Society of Human Reproduction and Embryology (ESHRE) recognize AMH as a useful marker for ovarian reserve assessment.

\section*{Why This Test Is Ordered}
\begin{itemize}
  \item Fertility evaluation: AMH is a marker of ovarian reserve that predicts the response to ovarian stimulation during in vitro fertilization (IVF), guiding the choice of stimulation protocol and informing patients about their reproductive potential
  \item Assessment of premature ovarian insufficiency (POI): a low AMH supports the diagnosis in women presenting with oligomenorrhea or amenorrhea before age 40
  \item Diagnosis and evaluation of polycystic ovary syndrome (PCOS): AMH is often elevated and is increasingly recognized as a useful adjunct to the Rotterdam criteria
  \item Counseling about the menopausal transition: AMH declines with age and can help predict the timing of menopause
  \item Prediction of the risk of ovarian hyperstimulation syndrome (OHSS): a high AMH before IVF identifies women at increased risk, enabling preventive measures
  \item Evaluation of the effect of ovarian surgery or gonadotoxic chemotherapy on ovarian reserve
  \item Investigation of primary amenorrhea and disorders of sex development in infants: in males, AMH is produced by Sertoli cells and is used to evaluate cryptorchidism, anorchia, and androgen insensitivity
  \item Monitoring of Sertoli cell function in males with suspected gonadal dysfunction
\end{itemize}

\section*{Primary vs Secondary Test}
This test may serve as a primary screening tool or as a secondary confirmatory test depending on the clinical context.

\section*{How This Test Is Performed}
For most patients, a health care professional will take a blood sample from a vein in your arm using a small needle. After the needle is inserted, a small amount of blood will be collected into a test tube or vial. You may feel a little sting when the needle goes in or out. This usually takes less than five minutes. For certain tests, a urine sample, stool sample, or other specimen may be required.

\section*{What Preparation Is Needed}
Preparation requirements vary by test. Some tests require fasting for 8 to 12 hours. Others have no special preparation. Your healthcare provider or laboratory will inform you of any specific preparation needed for this test.

\section*{Contraindications and Precautions}
\begin{itemize}
  \item No absolute contraindications. Factors that affect results: AMH levels are relatively stable throughout the menstrual cycle but may vary slightly between early follicular and luteal phases. Oral contraceptive pills may suppress AMH levels by 15-20%, potentially underestimating true ovarian reserve. AMH declines naturally with age, particularly after age 35, and becomes undetectable in most women after menopause. Endometriosis and ovarian surgery may reduce AMH levels. AMH should be interpreted alongside other ovarian reserve markers (FSH, estradiol, antral follicle count) for comprehensive assessment.
\end{itemize}
\section*{Risks}
Minimal risk. Slight pain or bruising at the venipuncture site.

\section*{What Happens After the Test}
Results are typically available within 1 to 3 days. AMH levels help guide fertility treatment planning: Low AMH (<1.0 ng/mL) suggests diminished ovarian reserve and may indicate need for more aggressive stimulation protocols or consideration of donor eggs. High AMH (>3.0-4.0 ng/mL) may suggest PCOS and increased risk of ovarian hyperstimulation syndrome (OHSS) during IVF stimulation. Normal AMH (1.0-3.0 ng/mL) suggests adequate ovarian reserve for conventional stimulation protocols.

\section*{Understanding Results}

\subsection*{Below Normal}
\begin{itemize}
  \item A low AMH (<1.0 ng/mL) indicates diminished ovarian reserve.
  \item This does not mean you cannot get pregnant, but it suggests: Lower response to ovarian stimulation medications, fewer eggs retrieved during IVF, earlier menopause than average.
  \item Causes of low AMH: natural age-related decline, premature ovarian insufficiency, prior ovarian surgery (endometrioma removal, cystectomy), chemotherapy or radiation therapy, autoimmune oophoritis.
  \item Low AMH does not predict natural fertility --- it specifically reflects the size of the antral follicle pool.
  \item Women with low AMH can still conceive spontaneously, but may have a shorter time to menopause.
\end{itemize}

\subsection*{Above Normal}
\begin{itemize}
  \item An elevated AMH (>3.0-4.0 ng/mL) is commonly associated with polycystic ovary syndrome (PCOS).
  \item PCOS is diagnosed by the Rotterdam criteria (presence of 2 of 3: oligo-anovulation, hyperandrogenism, polycystic ovaries on ultrasound).
  \item Other causes of elevated AMH: multiple small antral follicles, some ovarian granulosa cell tumors (though these typically produce very high levels >10 ng/mL).
  \item AMH alone does not diagnose PCOS --- clinical correlation with menstrual history, androgen levels, and ultrasound findings is required.
  \item High AMH predicts a good response to ovarian stimulation but also increased risk of OHSS.
\end{itemize}

\section*{Normal Results}
\begin{itemize}
  \item \textbf{Premenopausal women (reproductive age):} Varies by age. At age 25: approximately 2.5--6.8 ng/mL; at age 30: approximately 1.8--5.6 ng/mL; at age 35: approximately 1.2--4.0 ng/mL; at age 40: approximately 0.5--2.4 ng/mL; at age 45: approximately 0.1--1.2 ng/mL
  \item \textbf{Menopausal:} Undetectable or $<$ 0.08 ng/mL
  \item \textbf{Male (adult):} Approximately 0.4--5.9 ng/mL
\end{itemize}

\section*{Follow-up Tests}
\begin{itemize}
  \item \textbf{Antral Follicle Count (AFC):} Transvaginal ultrasound to directly count small follicles --- correlates with AMH and provides complementary information about ovarian reserve
  \item \textbf{Day 3 FSH and estradiol:} Traditional ovarian reserve markers; FSH $>$ 10 mIU/mL suggests diminished reserve. Used alongside AMH
  \item \textbf{Clomiphene Citrate Challenge Test (CCCT):} Historically used to assess ovarian reserve; largely replaced by AMH measurement
  \item \textbf{Inhibin B:} Another marker of ovarian reserve; less commonly measured than AMH
  \item \textbf{Thyroid function tests and prolactin:} If anovulation is present, to rule out other causes of infertility
\end{itemize}


\section*{References}
\bibliographystyle{unsrtnat}
\bibliography{references}

\clearpage
\section*{Regulatory Status and Authorization Date}
This chapter describes a test category, not a specific commercial product. FDA, CDSCO, EU IVDR, and MHRA approval or registration dates therefore cannot be assigned without the assay manufacturer, product name, instrument, and intended use. For laboratory-developed tests, an individual agency approval date may not exist. See the Preface for the interpretation of regulatory status.

\clearpage
